% --- IMPORTS ---
\documentclass[leqno]{article}
\usepackage{graphicx}
\usepackage{dsfont}
\usepackage{mathtools}
\usepackage{amssymb}
\usepackage{amsmath}
\usepackage{amsthm}
\usepackage{float}
\usepackage{ragged2e}
\usepackage{tensor}
\usepackage{fancyhdr}
\usepackage{empheq}
\usepackage[most]{tcolorbox}
\usepackage{enumitem}
\usepackage{dirtytalk}
\usepackage{marginnote}
\usepackage{microtype}
\usepackage[
  a4paper,
  left=0.8in,
  right=1in,
  top=1in,
  bottom=0.5in,
  headheight=14pt,
  headsep=0.4in
]{geometry}
\setlength{\parindent}{0cm}
% --- TikZ Packages ---
\usepackage{pgfplots}
\pgfplotsset{compat=1.18}
\usepgfplotslibrary{fillbetween, polar}
\usetikzlibrary{calc, positioning}
\usetikzlibrary{angles, quotes}
\usetikzlibrary{arrows.meta}
\usetikzlibrary{decorations.pathreplacing, decorations.markings}
\usetikzlibrary{patterns}
\usetikzlibrary{intersections}
% --- URL Packages ---
\usepackage[colorlinks=true,linkcolor=red,citecolor=magenta,urlcolor=black]{hyperref}%
\urlstyle{same}
% --- END OF IMPORTS ---

% --- CUSTOM COMMANDS ---
\RenewDocumentCommand{\marks}{o m}{%
  \IfValueTF{#1}%
    {\marginnote{\raggedright\textbf{[#2]}}[#1]}%
    {\marginnote{\raggedright\textbf{[#2]}}}%
}
\newcommand{\vspacerule}[2]{%
  \par\vspace*{#1}%
  \noindent\hrulefill\par
  \vspace*{#2}%
}
\definecolor{scarlet}{RGB}{205, 40, 75}
\newcommand{\questionitem}{%
  \item\phantomsection
  \addcontentsline{toc}{section}{Question \number\value{enumi}}%
}
\renewcommand{\contentsname}{Questions}
% --- END OF CUSTOM COMMANDS ---

% --- HEADER CONFIG ---
\pagestyle{fancy}
\fancyhead{}
\fancyfoot{}
\renewcommand{\headrulewidth}{0.9pt}
\lhead{Vector Product}
\chead{}
\rhead{\href{https://danieljsmith.org}{danieljsmith.org}}
% --- END OF HEADER CONFIG ---

\begin{document}
\vspace*{20pt}
\tableofcontents
\newpage
\begin{enumerate}[
  label=\textbf{\arabic*.},
  leftmargin=*,
  topsep=0pt,
  partopsep=0pt,
  parsep=0pt
]


% ---- Question 1 ----
\questionitem
% [Source: _QBT__Vector_Product, Q1]

\begin{enumerate}[label=\textbf{(\alph*)}]
  \item The vectors
  \[
  \mathbf{p}=\begin{pmatrix}1\\3\\2\end{pmatrix}
  \quad\text{and}\quad
  \mathbf{q}=\begin{pmatrix}2\\m\\-1\end{pmatrix}
  \]
  span a plane with normal vector
  \[
  \mathbf{n}=\begin{pmatrix}1\\-1\\1\end{pmatrix}
  \]
  Given that $\mathbf{q}\times\mathbf{p}$ is in the same direction as $\mathbf{n}$, find $m$ and the positive constant $\lambda$ such that
  \[
  \mathbf{q}\times\mathbf{p}=\lambda\mathbf{n}
  \]
  \marks[-20pt]{3}
  \item Using $\mathbf{q}\times\mathbf{p}$, determine the acute angle between $\mathbf{p}$ and $\mathbf{q}$. \marks{3} \\
\end{enumerate}

\newpage

% ---- Question 2 ----
\questionitem
% [Source: _QBT__Vector_Product, Q2]

You are given the fixed points $A(1,0,1)$ and $C(4,0,2)$, and the variable point $B(t,2,5-t)$, where $t$ is a real parameter. \\

\begin{enumerate}[label=\textbf{(\alph*)}]
  \item Using only the geometrical properties of the vector product, explain why the statement ``$\overrightarrow{AB} \times \overrightarrow{AC} = 0$" is false for all values of $t$. \marks{2} \\
  \item
  \begin{enumerate}[label=(\roman*)]
    \item Use the vector product to find an expression, in terms of $t$, for the area of triangle $ABC$. \marks{4}
    \item Hence determine the value of $t$ for which the area of triangle $ABC$ is a minimum. \marks{2} \\
  \end{enumerate}
\end{enumerate}

\newpage

% ---- Question 3 ----
\questionitem
% [Source: _QBT__Vector_Product, Q3]

With respect to a fixed origin $O$, the points $A$ and $B$ have coordinates $(1, 1, 0)$ and $(p, 3, 2)$ respectively, where $p$ is a constant. \\

For each of the following, determine the possible values of $p$ for which \\

\begin{enumerate}[label=\textbf{(\alph*)}]
  \item $\overrightarrow{AB}$ makes an angle of $60^\circ$ with the positive $y$-axis. \marks{3} \\
  \item $\overrightarrow{OA} \times \overrightarrow{AB}$ is parallel to
  \[
  \begin{pmatrix}2\\-2\\1\end{pmatrix}
  \]
  \marks[-20pt]{3}
  \item the area of triangle $OAB$ is $\tfrac{3}{2}$. \marks{3} \\
\end{enumerate}


\newpage


% ---- Question 4 ----
\questionitem
% [Source: _QBT__Vector_Product, Q4]

Points $P$, $Q$ and $R$ have position vectors $\mathbf{p}$, $\mathbf{q}$ and $\mathbf{r}$ respectively, relative to origin $O$. \\

It is given that $(\mathbf{q}-\mathbf{p}) \times (\mathbf{r}-\mathbf{q}) = 2\mathbf{p}$ and that $\lvert \mathbf{p} \rvert = 5$. \\

\begin{enumerate}[label=\textbf{(\alph*)}]
  \item Determine each of the following as either a single vector or a scalar quantity.
  \begin{enumerate}[label=(\roman*)]
    \item $(\mathbf{r}-\mathbf{q}) \times (\mathbf{q}-\mathbf{p})$ \marks{1}
    \item $\mathbf{p} \times \big((\mathbf{q}-\mathbf{p}) \times (\mathbf{r}-\mathbf{q})\big)$ \marks{2}
    \item $\mathbf{p} \cdot \big((\mathbf{q}-\mathbf{p}) \times (\mathbf{r}-\mathbf{q})\big)$ \marks{2} \\
  \end{enumerate}
  \item Describe a geometrical relationship between the points $O$, $P$, $Q$ and $R$ which can be deduced from
  \begin{enumerate}[label=(\roman*)]
    \item the statement $(\mathbf{q}-\mathbf{p}) \times (\mathbf{r}-\mathbf{q}) = 2\mathbf{p}$ \marks{1}
    \item the result of part (a)(iii) \marks{1} \\
  \end{enumerate}
\end{enumerate}

\newpage


% ---- Question 5 ----
\questionitem
% [Source: _QBT__Vector_Product, Q5]

\begin{enumerate}[label=\textbf{(\alph*)}]
  \item Let
  \[
  \mathbf{p}=\begin{pmatrix}1\\1\\2\end{pmatrix}
  \quad\text{and}\quad
  \mathbf{q}=\begin{pmatrix}x\\0\\1\end{pmatrix}
  \]
  The area of the parallelogram formed by $\mathbf{p}$ and $\mathbf{q}$ is $\sqrt{2}$.\\
  It is also given that $\mathbf{p}\times\mathbf{q}$ is perpendicular to
  \[
  \begin{pmatrix}1\\3\\m\end{pmatrix}
  \]
  Find $x$ and $m$. \marks{4} \\
  \item Hence, using $|\mathbf{p}\times\mathbf{q}|$, determine the acute angle between $\mathbf{p}$ and $\mathbf{q}$. \marks{3} \\
\end{enumerate}
\newpage


% ---- Question 6 ----
\questionitem
% [Source: _QBT__Vector_Product, Q6]

The point $P$ has coordinates $P(2,-1,1)$. The vector $\overrightarrow{PQ}=(1,2,0)$ and the midpoint of $PR$ is $M(1,0,2)$. \\

\begin{enumerate}[label=\textbf{(\alph*)}]
  \item Use a vector product to show that the area of triangle $PQR$ is $\sqrt{14}$. \marks{4} \\
  \item Find a vector equation of the plane containing $P$, $Q$ and $R$ in the form $\mathbf{r} \cdot \mathbf{n} = k$. \marks{1} \\
  \item Hence find the exact distance of the plane from the origin. \marks{1} \\
\end{enumerate}

\newpage

% ---- Question 7 ----
\questionitem
% [Source: _QBT__Vector_Product, Q7]

The points $P$, $Q$ and $R$ have position vectors $\mathbf{p}=2\mathbf{i}-2\mathbf{k}$, $\mathbf{q}=3\mathbf{j}+3\mathbf{k}$ and $\mathbf{r}=t\mathbf{i}-\mathbf{j}+\mathbf{k}$ respectively, relative to the origin $O$. \\

Determine the value(s) of $t$ in each of the following cases. \\

\begin{enumerate}[label=\textbf{(\alph*)}]
  \item The line $OR$ is parallel to a normal to the plane $OPQ$. \marks{2} \\
  \item The volume of tetrahedron $OPQR$ is $9$. \marks{4} \\
\end{enumerate}

\newpage

% ---- Question 8 ----
\questionitem
% [Source: _QBT__Vector_Product, Q8]

The points $A$, $B$ and $C$ lie on the plane $\Pi$. \\

The position vectors of $A$ and $B$ are $\mathbf{a}=2\mathbf{i}-\mathbf{j}+3\mathbf{k}$ and $\mathbf{b}=5\mathbf{i}+\mathbf{j}$ respectively. \\

The centroid $G$ of triangle $ABC$ has position vector $4\mathbf{i}+\mathbf{j}+2\mathbf{k}$. \\

\begin{enumerate}[label=\textbf{(\alph*)}]
  \item Find $\overrightarrow{AB} \times \overrightarrow{AC}$. \marks{4} \\
  \item Find an equation for $\Pi$ in the form $\mathbf{r} \cdot \mathbf{n} = p$. \marks{2} \\
\end{enumerate}
The point $D$ has position vector $\mathbf{i}+5\mathbf{j}+4\mathbf{k}$. \\
\begin{enumerate}[label=\textbf{(\alph*)}, resume]
  \item Determine the volume of the tetrahedron $ABCD$. \marks{4} \\
\end{enumerate}

\newpage

% ---- Question 9 ----
\questionitem
% [Source: _QBT__Vector_Product, Q9]

With respect to a fixed origin $O$ the point $A$ has coordinates $(9,10,-1)$ and the line $\ell$ is the intersection of the planes $x-z=-2$ and $x+2y=5$. \\

The points $B$ and $C$ lie on $\ell$ such that $AB=AC=15$. \\

Given that $A$ does not lie on $\ell$ and that the $z$ coordinate of $B$ is negative \\

\begin{enumerate}[label=\textbf{(\alph*)}]
  \item determine the coordinates of $B$ and the coordinates of $C$. \marks{4} \\
  \item Hence determine a Cartesian equation of the plane containing the points $A$, $B$ and $C$. \marks{3} \\
\end{enumerate}
The point $D$ has coordinates $(8,4,\alpha)$, where $\alpha$ is a constant. \\

Given that the volume of the tetrahedron $ABCD$ is $156$ \\

\begin{enumerate}[label=\textbf{(\alph*)}, resume]
  \item determine the possible values of $\alpha$. \marks{4} \\
\end{enumerate}
Given that $\alpha>0$ \\

\begin{enumerate}[label=\textbf{(\alph*)}, resume]
  \item determine the shortest distance between the line $\ell$ and the line passing through the points $A$ and $D$, giving your answer to $2$ significant figures. \marks{4} \\
\end{enumerate}

\newpage

% ---- Question 10 ----
\questionitem
% [Source: _QBT__Vector_Product, Q10]

The lines $L_1$, $L_2$ and $L_3$ are given by
\begin{align*}
L_1: &\quad \mathbf{r}=\begin{pmatrix}8\\3\\1\end{pmatrix}+\lambda\begin{pmatrix}3\\2\\-1\end{pmatrix}\\[6pt]
L_2: &\quad x=9-4\mu,\ y=-2+3\mu,\ z=-1+3\mu\\[6pt]
L_3: &\quad \dfrac{x-3}{1}=\dfrac{y+6}{-5}=\dfrac{z-1}{-2}
\end{align*}
The pairwise intersections of these lines are the vertices of a triangle. \\

Find the area of this triangle. \marks{9}


\end{enumerate}
\end{document}
